% Cote 119: offprints, reports, courses, CV and letters (typed and printed part): archive cote 119, site page 61 (PDF page 62).
% offprint. Draft compared: none.
% AI-transcribed by Claude Opus 5.5 (claude-opus-5-5, adaptive thinking, effort high), Message Batches API, batch pilotnext (msgbatch_01A9sdu8SYq23bEtqX9tUncQ), 2026-11-13; not yet checked by eye.
\scanpage{119/61}
\mspage{C. R. Acad. Sc. Paris, t. 273, p. 1144-1147 (8 décembre 1971) \quad Série A}

GÉOMÉTRIE ALGÉBRIQUE. --- Polarisations des catégories tannakiennes :
cas homogène. Note (*) de M. Neantro Saavedra Rivano, transmise par
M. Henri Cartan.

Dans cette Note, on définit et on étudie des « structures de positivité » sur des
$\otimes$-catégories de la forme $\mathit{Rep}_0\,(\mathrm{G})$ pour G un $\mathbf{R}$-groupe algébrique affine, et plus
généralement sur des catégories tannakiennes sur $\mathbf{R}$ [voir (${}^1$), 2.1].
On donne un sous-corps K de $\mathbf{R}$ (qui sera $\mathbf{R}$ lui-même à partir du 2.3) et une
catégorie tannakienne $\mathcal{C}$ sur K [par exemple, $\mathit{Rep}_0\,(\mathrm{G})$ pour un K-groupe affine G]
dont les objets seront notés V, W, \dots. Enfin, Z désigne le K-groupe affine abélien,
centre du lien de $\mathcal{C}$; on a un isomorphisme canonique
\[ \mathrm{Z} \xrightarrow{\sim} \mathrm{Aut}^{\otimes}(\mathrm{id}_{\mathcal{C}}). \]

1. Formes de Weil. --- 1.1. Si $\varphi : \mathrm{V}\otimes\mathrm{V} \to 1$ est une forme bilinéaire
non-dégénérée sur V à valeurs dans l'objet unité de $\mathcal{C}$, on définit la parité $\varepsilon_{\varphi}$
de $\varphi$ comme le seul endomorphisme de V vérifiant
\[ (1.1.1) \qquad \varphi\circ(\mathrm{id}_{\mathrm{V}}\otimes\varepsilon_{\varphi}) = \varphi\circ\sigma_{\mathrm{V}}, \]
où $\sigma_{\mathrm{V}} : \mathrm{V}\otimes\mathrm{V} \xrightarrow{\sim} \mathrm{V}\otimes\mathrm{V}$ est la symétrie canonique donnée par la contrainte
de commutativité sur la loi $\otimes$ de $\mathcal{C}$. On définit aussi, si $u\in\mathrm{End}\,(\mathrm{V})$,
le transposé $u^{\varphi}$ de $u$ relativement à la forme $\varphi$; on a
\[ (1.1.2) \qquad \varphi\circ(u\otimes\mathrm{id}_{\mathrm{V}}) = \varphi\circ(\mathrm{id}_{\mathrm{V}}\otimes u^{\varphi}); \]
\[ (1.1.3) \qquad (u^{\varphi})^{\varphi} = \varepsilon_{\varphi}\, u\, \varepsilon_{\varphi}^{-1}. \]

Définition 1.2. --- Une forme bilinéaire non-dégénérée $\varphi : \mathrm{V}\otimes\mathrm{V} \to 1$
est appelée une forme de Weil si $\varepsilon_{\varphi}$ est dans le centre de End (V) et si,
pour $u\in\mathrm{End}\,(\mathrm{V})$, $u\neq 0$, on a $\mathrm{Tr}\,(uu^{\varphi}) > 0$. Deux formes de Weil
$\varphi : \mathrm{V}\otimes\mathrm{V} \to 1$, $\psi : \mathrm{W}\otimes\mathrm{W} \to 1$ sont compatibles si $\varphi\oplus\psi$ est une forme
de Weil sur $\mathrm{V}\oplus\mathrm{W}$.
1.3. Si $\varepsilon$ est un automorphisme de V central dans End (V), la relation
de compatibilité pour des formes de Weil sur V de parité $\varepsilon$ est une relation
d'équivalence; l'ensemble quotient sera noté $w_{\varepsilon}\,(\mathrm{V})$. Si $\mathrm{V}\ \ = \mathrm{V}\otimes_{\mathrm{K}}\mathbf{R}$ est
l'objet déduit de V dans la catégorie tannakienne $\mathcal{C}_{\mathbf{R}}$ sur $\mathbf{R}$ obtenue de $\mathcal{C}$
par extension des scalaires de K à $\mathbf{R}$ (${}^2$), il résulte aussitôt de la densité
de K dans $\mathbf{R}$ qu'on a une bijection canonique
\[ (1.3.1) \qquad w_{\varepsilon}\,(\mathrm{V}) \simeq w_{\varepsilon}\,(\mathrm{V}_{\mathbf{R}}). \]

Si V est un objet semi-simple de $\mathcal{C}$ et $\varphi$ une forme de Weil sur V, pour tout
sous-objet V$'$ de V, la restriction $\varphi_{\mathrm{V}'}$ de $\varphi$ à V$'$ est une forme de Weil sur V$'$,
compatible avec $\varphi$. De plus, l'ensemble $w_{\varepsilon}\,(\mathrm{V})$ $(\varepsilon = \varepsilon_{\varphi})$ est fini et a
$2^{r}$ éléments, où $r$ est le nombre de composantes isotypiques de V.
Exemple 1.4. --- Soient G un $\mathbf{R}$-groupe affine, $\mathrm{C}\in\mathrm{G}\,(\mathbf{R})$ tel que $\mathrm{C}^{2}\in\mathrm{Z}\,(\mathbf{R})$,
où Z est le centre de G. On pose $\mathcal{C} = \mathit{Rep}_0\,(\mathrm{G})$. Si V est un objet de $\mathcal{C}$,

\mspage{\hand{61}}
